\documentclass[11pt]{article}
\usepackage[a4paper,margin=23mm]{geometry}
\usepackage{amsmath,amssymb}
\usepackage[hidelinks]{hyperref}
\title{RM0016: printed off-ramp queue consistency check}
\author{Review note, 14 September 2026}
\date{}
\begin{document}
\maketitle
\noindent\textbf{Scope.} Zhang et al., DOI 10.1109/MITS.2023.3323029, supplied final PDF p.~7 (printed p.~46), Eqs.~(7)--(8). The equations were checked against a rendered page, not only extracted text. This note assesses the printed recurrence, not an unavailable implementation.

\paragraph{Notation.} Write $Q_t=N^{\mathrm{OFF}}_{j,t}$ for off-ramp queue mass, $S=S^{\mathrm{OFF}}_j$ for storage, $d_{t+1}=q_{n-1,j,t+1}$ for arriving demand rate, $c=e_j^{\max}$ for downstream capacity rate, and $T>0$ for step duration. The printed equations become
\begin{align}
e_{t+1}&=\min\left\{d_{t+1},\frac{S-Q_t}{T}\right\},\\
Q_{t+1}&=\max\{Q_t+T(e_{t+1}-c),0\}.
\end{align}
Assume fixed $S\geq0$, $c\geq0$, nonnegative demand and $0\leq Q_t\leq S$. These are the ordinary feasible-state conditions, not an additional control constraint.

\paragraph{One-step bound.} The first equation implies
\[
T e_{t+1}\leq S-Q_t.
\]
Consequently,
\[
Q_{t+1}\leq\max\{S-Tc,0\}\leq S.
\]
Nonnegativity follows from the second equation. By induction, a feasible initial queue remains at or below storage. In particular, a trigger defined by $Q_t>S$ cannot activate through this recurrence from a feasible initial state.

\paragraph{Why clarification is needed.} The article describes a lane-drop trigger when the off-ramp queue exceeds storage, and reports relative queues above one (e.g., Fig.~9). Those claims cannot be generated by the literal recurrence under the conditions above. Soft overflow penalties do not change this bound. A different plant update, different queue definition, changing storage, initially overflowing state, or printing error would need to be specified explicitly.

\paragraph{Review consequence.} Do not silently replace the equations, infer misconduct, or use this discrepancy as a proof that ramp metering is infeasible. It blocks unqualified literal use of this source. Its conceptual role and the wider B1 feasibility question must be judged separately.
\end{document}
